% Q4V3: TL System — Short-Circuit Load, Mismatched Source
% Source (Vg, Rg≠Z0) → TL (Z0, ℓ) → Short circuit
\documentclass[border=10pt]{standalone}
\usepackage[american voltages, american currents]{circuitikz}
\usepackage{amsmath}
\usetikzlibrary{decorations.pathreplacing}
\renewcommand{\familydefault}{\sfdefault}

\begin{document}
\begin{circuitikz}[line width=0.8pt, every node/.style={font=\sffamily}]

% === Source side ===
\draw (0,0) to[V, v=$V_g$] (0,3);
\draw (0,3) to[R, l=$R_g$] (2.5,3);

% === Transmission line (two parallel lines) ===
\draw[double, double distance=6pt, line width=0.8pt] (2.5,3) -- (7.5,3);
\draw[double, double distance=6pt, line width=0.8pt] (2.5,0) -- (7.5,0);

% TL labels
\node[above, font=\sffamily] at (5,3.4) {$Z_0$, $\ell$};
\node[below, font=\sffamily\small] at (5,-0.4) {$v_p$};

% === Short circuit load (wire connecting top to bottom) ===
\draw[line width=1.2pt] (7.5,3) -- (7.5,0);

% Short circuit label
\node[right, font=\sffamily\small] at (7.8,1.5) {Short};
\node[right, font=\sffamily\small] at (7.8,1.0) {circuit};

% === Bottom wire (source side) ===
\draw (0,0) -- (2.5,0);

% === Junction dots ===
\fill (2.5,3) circle (2pt);
\fill (2.5,0) circle (2pt);
\fill (7.5,3) circle (2pt);
\fill (7.5,0) circle (2pt);

% === Position markers ===
\draw[dashed, gray] (2.5,-0.8) -- (2.5,0);
\node[below, font=\sffamily\small, gray!70!black] at (2.5,-0.9) {$z = 0$};
\draw[dashed, gray] (7.5,-0.8) -- (7.5,0);
\node[below, font=\sffamily\small, gray!70!black] at (7.5,-0.9) {$z = \ell$};

% === Source label (below Vg, clear of polarity marks) ===
\node[below, font=\sffamily\small] at (0,-0.4) {Source};

\end{circuitikz}
\end{document}
